\section{Proofs}
\label{app:proofs}

\subsection{\Cref{prop:lmmse}}
Write $\lmsg^{\mathrm{p}}_i$ for the forward message at site $i$ after projection and
$\ell_i$ for the Gaussian unary likelihood. We show by induction that every projected
message equals the corresponding exact message of the covariance-matched Gaussian chain.
The initial message is Gaussian or projected to the Gaussian with the same first two
moments, which is the Gaussian chain's initial message. If $\lmsg^{\mathrm{p}}_{i-1}$ is
that chain's message, with mean $m$ and variance $v$, multiplying by $\ell_{i-1}$ gives
a Gaussian with the precision-weighted moments, identical in both chains; the transition
integral against $a_i=\corr a_{i-1}+\varepsilon_i$ is a convolution whose mean and
variance are $(\corr m,\corr^2v+\ivar)$ for every innovation with moments $(0,\ivar)$, so
$\Pgauss$ returns the Gaussian chain's next message. The backward sweep runs the same
argument with the inverted map $(m,v)\mapsto(m/\corr,(v+\ivar)/\corr^2)$, which is where
$\corr\neq0$ is used, starting from the terminal likelihood $\ell_{\nsites}$ (the flat
terminal message is not normalisable and contributes no factor). Every projected belief
is therefore the Gaussian chain's belief, whose posterior mean is linear in $\vx$ and
coincides with the LMMSE estimator $e^{-t}\Sig_0\Sig_t^{-1}\vx$. Since the recursion
never read $\innov$ beyond $(0,\ivar)$, the estimator is the same for every such law.
The algebraic steps (the two moment maps as integrals, the covariance identities, and
the optimality of the normal-equation gain in arbitrary dimension) are machine-checked
in Lean~4 (\codefile{lean/README.md}); the projection step itself is not.

\subsection{\Cref{thm:hermite}}
For $z\sim\normal{0}{1}$ and real $a$, $\alpha$ with $\beta=\sqrt{1-\alpha^2}$, the
generating function $e^{sx-s^2/2}=\sum_n\Hen{n}(x)s^n/n!$ evaluated at
$x=\alpha a+\beta z$ and averaged over $z$ gives
$e^{s\alpha a-s^2\alpha^2/2}=\sum_n\Hen{n}(a)(\alpha s)^n/n!$, hence
$\E_z\Hen{n}(\alpha a+\beta z)=\alpha^n\Hen{n}(a)$. The coordinates of $\vz$ are
independent given $\va$, so the product form holds for $H_\nu$; integrating over
$\va\sim\Pz$ requires only $\E|\va|^{|\nu|}<\infty$. The cases $n\le4$ are verified in
Lean as polynomial identities under the first four Gaussian moments.

\subsection{\Cref{thm:decay}}
\label{app:proof-decay}

\paragraph{Exact identities.}
With $\alpha=e^{-t}$ and $\Dt=1-\alpha^2$, Bayes' rule gives
$\Pt(\vx)=\gamma_{\Dt}(\vx)\,\E_{\Pz}\exp\bigl(\alpha\vx\cdot\va/\Dt-\alpha^2|\va|^2/2\Dt\bigr)$,
where $\gamma_{\Dt}$ is the $\normal{0}{\Dt\Id}$ density. Hence
$\log\Pt(\vx)=\log\gamma_{\Dt}(\vx)+\Lambda(s,b)$ with $s=\alpha\vx/\Dt$, $b=\alpha^2/\Dt$
and $\Lambda(s,b)=\log\E_{\Pz}\exp(s\cdot\va-b|\va|^2/2)$. For $b>0$ the tilted measure
is integrable for every $s$, so $\Lambda$ is real-analytic in $(s,b)$ on
$\Reals^{\nsites}\times(0,\infty)$, and its Taylor coefficients at the origin are the
joint cumulants of the linear form $s\cdot\va$ and the quadratic form $|\va|^2/2$;
$\E[\va\mid\vx]=\nabla_s\Lambda$ and $\score=-\vx/\Dt+(\alpha/\Dt)\nabla_s\Lambda$.

\paragraph{Grading.}
Expand $\Lambda$ in $(s,b)$ at the origin and grade a monomial $s^{j}b^{k}$ by $j+2k$;
since $s=O(\alpha)$ and $b=O(\alpha^2)$, the grade is the order in $\alpha$. By the
symmetry of $\Pz$ under $\va\mapsto-\va$ only even grades occur:
$\Lambda=\Lambda_2+\Lambda_4+R_6$, with
$\Lambda_2=\tfrac12s^{\!\top}\Sig_0s-\tfrac b2\operatorname{tr}\Sig_0$ and
\[
  \Lambda_4=\tfrac1{24}\kfour[s,s,s,s]
   -\tfrac b4\sum_i\operatorname{cum}(s\!\cdot\!\va,s\!\cdot\!\va,a_i,a_i)
   +\tfrac{b^2}8\sum_{ij}\operatorname{cum}(a_i,a_i,a_j,a_j),
\]
where $\operatorname{cum}$ is the joint cumulant under $\Pz$. Every fourth-order joint
cumulant of linear forms in $\va$ equals its value under the covariance twin plus a
linear contraction of $\kfour$, so $\Lambda_4-\Lambda_4^{\mathrm{G}}$ is linear in
$\kfour$ and $\Lambda_2=\Lambda_2^{\mathrm{G}}$. The remainder $R_6$ is the Taylor
remainder of grade six; on $\{|\alpha\vx|\le1\}$ it is bounded by
$C\alpha^6(1+|\vx|)^6$ with $C$ depending on the moments of $\Pz$ up to order eight.

\paragraph{Leading term of the score gap.}
$\score-\Sgauss=(\alpha/\Dt)\nabla_s(\Lambda-\Lambda^{\mathrm{G}})
=(\alpha/\Dt)\bigl[\nabla_s(\Lambda_4-\Lambda_4^{\mathrm{G}})+\nabla_sR_6-\nabla_sR_6^{\mathrm{G}}\bigr]$.
Substituting $s=\alpha\vx/\Dt$, $b=\alpha^2/\Dt$ and expanding in $\alpha$, the
grade-four term is
$\alpha^4\,\nabla_{\vx}\bigl[\tfrac1{24}\sum_{ijkl}\kfour[ijkl]\,{:}x_ix_jx_kx_l{:}\bigr]+O(\alpha^6(1+|\vx|)^5)$,
where ${:}\cdot{:}$ is the Wick (Hermite) product: the three contributions of $\Lambda_4$
combine into the Hermite form because $b=\alpha^2(1+O(\alpha^2))$ and the cumulant
contractions with $\delta_{ii}$ are exactly the terms that convert monomials into Hermite
polynomials. The $m$-th component is $\alpha^4\,\tfrac16\sum_{jkl}\kfour[mjkl]\,{:}x_jx_kx_l{:}$.

\paragraph{Second moment.}
Under $\gamma$, Wick's theorem gives
$\E_\gamma[{:}x_jx_kx_l{:}\,{:}x_{j'}x_{k'}x_{l'}{:}]$ equal to the number of perfect
matchings of $\{j,k,l\}$ with $\{j',k',l'\}$ respecting indices; contracted with the
symmetric tensor this yields
$\E_\gamma\bigl[(\tfrac16\sum_{jkl}\kfour[mjkl]{:}x_jx_kx_l{:})^2\bigr]
=\tfrac{3!}{36}\sum_{jkl}\kfour[mjkl]^2$, and summing over $m$,
$\E_\gamma\lVert\text{leading term}\rVert^2=\tfrac16\Frob{\kfour}^2\alpha^8$. Replacing
$\E_\gamma$ by $\E_{\Pt}$ costs a factor $1+O(\alpha^2\lmax)$: on the bulk,
$\Pt/\gamma=1+\alpha^2\tfrac12\langle\Sig_0-\Id,{:}\vx\vx^{\!\top}{:}\rangle+O(\alpha^4\lmax^2(1+|\vx|)^4)$,
and the degree-2 correction has zero $\gamma$-mean against the square of a degree-3
polynomial up to terms of relative size $\alpha^2\lmax$. The cross term between the
leading term and the grade-six remainder is $O(\alpha^{10}\lmax)$ by Cauchy--Schwarz and
the moment bound.

\paragraph{Tail.}
On $\{|\alpha\vx|>1\}$ the expansion is not used. There $|\score-\Sgauss|\le
2|\vx|/\Dt+\alpha(\E[|\va|\,|\vx]+|\E^{\mathrm G}[\va|\vx]|)/\Dt\le C(1+|\vx|)/\Dt$
because the tilted mean is bounded by a polynomial in $|\vx|$ under the moment
condition, while $\Pt(|\alpha\vx|>1)\le\exp(-c/\alpha^2)$ since $\vx_t$ is a Gaussian
convolution with variance at most one per coordinate. The tail contribution is therefore
$O(\exp(-c'e^{2t}))$, below every power of $\alpha$. Together,
$\E_{\Pt}\lVert\score-\Sgauss\rVert^2=\tfrac16\Frob{\kfour}^2e^{-8t}(1+O(e^{-2t}\lmax))$,
which is \cref{eq:scoregap}.

\paragraph{Information.}
$\partial_\theta\log\Pt=\partial_\theta\Lambda$. If $\partial_\theta\Sig_0\neq0$ the
grade-two term $\tfrac{\alpha^2}2\langle\partial_\theta\Sig_0,{:}\vx\vx^{\!\top}{:}\rangle$
leads, and $\E_\gamma$ of its square is
$\tfrac{\alpha^4}4\cdot2\Frob{\partial_\theta\Sig_0}^2$ by Wick, giving \cref{eq:info2}.
If $\partial_\theta\Sig_0=0$ the grade-four term
$\tfrac{\alpha^4}{24}\sum\partial_\theta\kfour[ijkl]{:}x_ix_jx_kx_l{:}$ leads, with
$\gamma$-second moment $\tfrac{4!}{24^2}\Frob{\partial_\theta\kfour}^2\alpha^8$, giving
\cref{eq:info4}; the corrections and the tail are handled as above. Since
${:}\vx\vx^{\!\top}{:}$ and the degree-four Hermite polynomials are $\gamma$-orthogonal,
the leading cross-information between a degree-2 and a degree-4 coordinate vanishes, so
the leading term of the efficient information of the degree-4 coordinate equals that of
its raw information. For the Gaussian chain $\vx_t\sim\normal{0}{\Sig_t(\theta)}$ exactly
and the stated trace formula is the Gaussian Fisher information; its limit
$e^{4t}I_t\to\tfrac12\Frob{\partial_\theta\Sig_0}^2$ follows from $\Sig_t\to\Id$ and
$\partial_\theta\Sig_t=\alpha^2\partial_\theta\Sig_0$, and is machine-checked in Lean
(\codefile{lean/README.md}, target~3).

\paragraph{The crossover.}
The expansion of $\Pt/\gamma$ used above is an expansion in $\alpha^2\Sig_0$, i.e.\ in
$\alpha^2\lmax(\Sig_0)$ in operator norm; for the chain $\lmax\le(1+\corr)/(1-\corr)$ by
the Toeplitz bound $\lmax\le\sum_k\corr^{|k|}$. This is the source of the factor
$\lmax$ in every remainder and of $\tstar=\tfrac12\log\lmax$; for the Gaussian chain the
exact formula shows the same scale directly, since $\Sig_t^{-1}$ is within
$\alpha^2\lmax$ of $\Id$ in operator norm.
